\section{Experiments}

\subsection{Datasets}

We evaluate DVCL on four heterogeneous graph benchmarks: ACM, DBLP, AMiner, and IMDB. Each dataset contains multiple node types and relation types, and the task is semi-supervised node classification on a target node type. Following the HSeCo setting, we use dataset-specific meta-paths to construct semantic neighborhoods for target nodes. Table~\ref{tab:dataset_settings} summarizes the target node type and meta-paths used in each dataset.

\begin{table}[t]
\centering
\caption{Dataset settings.}
\label{tab:dataset_settings}
\begin{tabular}{c|c|c}
\hline
Dataset & Target node type & Meta-paths \\
\hline
ACM & Paper & P-A-P, P-F-P \\
DBLP & Author & A-P-A, A-P-C-P-A, A-P-T-P-A \\
AMiner & Paper & P-A-P, P-R-P \\
IMDB & Movie & M-D-M, M-A-M \\
\hline
\end{tabular}
\end{table}

The current completed main experiments focus on ACM, where both global structural attacks and target attacks have been evaluated with five random seeds. DBLP results are available as single-seed supporting evidence. AMiner and IMDB are kept in the benchmark protocol and will be completed in the next experimental round.

\subsection{Baselines}

The primary baseline is HSeCo, because DVCL is built on top of HSeCo's semantic topology learning pipeline and is designed to test whether adding a feature-induced view and cross-view contrastive alignment improves robustness. We also position HSeCo against other related heterogeneous graph and robust graph learning methods, including HAN~\cite{wang2019han}, HGT~\cite{hu2020hgt}, MAGNN~\cite{fu2020magnn}, HeCo~\cite{wang2021heco}, RoHe~\cite{rohe}, HeteGuard~\cite{heteguard}, GCN-Jaccard~\cite{wu2019gcnjaccard}, RGCN~\cite{zhu2019rgcn}, Pro-GNN~\cite{jin2020prognn}, GNNGuard~\cite{zhang2020gnnguard}, and SimP-GCN~\cite{jin2021simpgcn}.

In the current draft, the reported numeric comparisons use HSeCo as the direct baseline. DVCL variants are used only for ablation and supplementary analysis, not as external baselines.

\subsection{Attack Settings}

For global attacks, we evaluate PRBCD and HetePRBCD with perturbation rates of $5\%$, $10\%$, $15\%$, $20\%$, and $25\%$. PRBCD performs global structure perturbation, while HetePRBCD accounts for heterogeneous relation types during perturbation. Clean performance without attack is also reported.

For target attacks, we use an FGSM-style heterogeneous graph target attack on ACM. For each seed, 500 target paper nodes are sampled from the test split. The attack budget is set to $\Delta=3$ in the main multi-seed target-attack table, and the attack perturbs task-relevant P-A edges. The target-attack metric is computed on the selected target nodes, not on the full test split.

\subsection{Implementation Details}

All models are trained with the same data splits and attack artifacts for fair comparison. Unless otherwise specified, each ACM result is averaged over five random seeds, namely 1, 2, 3, 4, and 5. The maximum number of training epochs is 200, and early stopping is used. We use Adam as the optimizer.

The main DVCL setting uses the original feature KNN view: raw target-node features, directed KNN, and $k=20$. The topology view follows the HSeCo semantic topology construction. The two view-specific representations are concatenated for classification, and the cross-view contrastive loss is applied between topology-view and feature-view embeddings. We use $\lambda_{\mathrm{HAN}}=1.0$. The main tuned DVCL configuration uses $\lambda_{\mathrm{DVCL}}=0.5$; the earlier $\lambda_{\mathrm{DVCL}}=0.1$ setting is retained in Table~\ref{tab:lambda_compare} for comparison.

\subsection{Overall Performance}

Table~\ref{tab:overall_global} reports the main ACM global-attack results. HSeCo is the direct baseline. DVCL uses the raw directed KNN feature view with $k=20$ and $\lambda_{\mathrm{DVCL}}=0.5$.

\begin{table}[t]
\centering
\caption{ACM global attack performance. Results are Micro-F1 mean $\pm$ standard deviation over five seeds.}
\label{tab:overall_global}
\begin{tabular}{c|c|c|c|c}
\hline
Attack & Rate & HSeCo & DVCL & $\Delta$ pp \\
\hline
Clean & 0\% & $0.8679\pm0.0129$ & $0.8770\pm0.0071$ & +0.91 \\
PRBCD & 5\% & $0.8491\pm0.0086$ & $0.8732\pm0.0092$ & +2.41 \\
PRBCD & 10\% & $0.8519\pm0.0080$ & $0.8718\pm0.0077$ & +1.99 \\
PRBCD & 15\% & $0.8484\pm0.0091$ & $0.8674\pm0.0060$ & +1.90 \\
PRBCD & 20\% & $0.8530\pm0.0092$ & $0.8660\pm0.0146$ & +1.30 \\
PRBCD & 25\% & $0.8504\pm0.0092$ & $0.8659\pm0.0082$ & +1.55 \\
HetePRBCD & 5\% & $0.8530\pm0.0121$ & $0.8762\pm0.0027$ & +2.33 \\
HetePRBCD & 10\% & $0.8589\pm0.0126$ & $0.8737\pm0.0087$ & +1.48 \\
HetePRBCD & 15\% & $0.8546\pm0.0070$ & $0.8726\pm0.0046$ & +1.79 \\
HetePRBCD & 20\% & $0.8548\pm0.0124$ & $0.8700\pm0.0108$ & +1.53 \\
HetePRBCD & 25\% & $0.8511\pm0.0092$ & $0.8703\pm0.0063$ & +1.92 \\
\hline
\end{tabular}
\end{table}

DVCL outperforms HSeCo on all clean and attacked ACM settings. The average Micro-F1 over the ten attack conditions is 0.8525 for HSeCo and 0.8707 for DVCL, corresponding to an average gain of 1.82 percentage points. The improvement is consistent across both PRBCD and HetePRBCD, showing that the feature-induced view and cross-view alignment improve robustness beyond HSeCo's topology-only semantic defense.

Table~\ref{tab:overall_macro} reports the corresponding Macro-F1 comparison. The Macro-F1 trend is consistent with Micro-F1, indicating that the improvement is not limited to dominant classes.

\begin{table}[t]
\centering
\caption{ACM global attack Macro-F1. Results are mean $\pm$ standard deviation over five seeds.}
\label{tab:overall_macro}
\begin{tabular}{c|c|c|c|c}
\hline
Attack & Rate & HSeCo & DVCL & $\Delta$ pp \\
\hline
Clean & 0\% & $0.8654\pm0.0129$ & $0.8739\pm0.0074$ & +0.85 \\
PRBCD & 5\% & $0.8450\pm0.0090$ & $0.8709\pm0.0085$ & +2.59 \\
PRBCD & 10\% & $0.8475\pm0.0085$ & $0.8680\pm0.0076$ & +2.05 \\
PRBCD & 15\% & $0.8446\pm0.0099$ & $0.8639\pm0.0059$ & +1.93 \\
PRBCD & 20\% & $0.8497\pm0.0100$ & $0.8641\pm0.0148$ & +1.44 \\
PRBCD & 25\% & $0.8473\pm0.0093$ & $0.8620\pm0.0108$ & +1.47 \\
HetePRBCD & 5\% & $0.8490\pm0.0125$ & $0.8738\pm0.0041$ & +2.48 \\
HetePRBCD & 10\% & $0.8557\pm0.0127$ & $0.8706\pm0.0088$ & +1.49 \\
HetePRBCD & 15\% & $0.8516\pm0.0059$ & $0.8704\pm0.0042$ & +1.88 \\
HetePRBCD & 20\% & $0.8511\pm0.0124$ & $0.8677\pm0.0109$ & +1.66 \\
HetePRBCD & 25\% & $0.8467\pm0.0090$ & $0.8668\pm0.0070$ & +2.01 \\
\hline
\end{tabular}
\end{table}

\subsection{Target Attack Results}

Table~\ref{tab:target_attack} reports ACM target-attack results under budget $\Delta=3$. HSeCo drops by 10.16 percentage points in target Micro-F1, while DVCL drops by only 1.16 percentage points. Under the same target attack setting, DVCL improves attacked target Micro-F1 by 10.00 percentage points over HSeCo.

\begin{table}[t]
\centering
\caption{ACM FGSM-style target attack under budget $\Delta=3$. Results are averaged over five seeds.}
\label{tab:target_attack}
\begin{tabular}{c|c|c|c|c}
\hline
Model & Clean Target Micro-F1 & Attacked Target Micro-F1 & Drop & Attacked Target Macro-F1 \\
\hline
HSeCo & 0.8624 & 0.7608 & -10.16 pp & 0.7631 \\
DVCL & 0.8724 & 0.8608 & -1.16 pp & 0.8585 \\
\hline
\end{tabular}
\end{table}

This result shows the largest observed robustness gap between HSeCo and DVCL. Since target attacks directly perturb P-A edges around selected target nodes, the topology view can be strongly affected. DVCL is more stable because the raw feature KNN view does not directly depend on the attacked edges and can compensate for corrupted semantic topology.

\subsection{Ablation Study}

We first compare the original DVCL setting with the tuned contrastive setting. Table~\ref{tab:lambda_compare} shows that increasing $\lambda_{\mathrm{DVCL}}$ from 0.1 to 0.5 improves all attacked global settings while keeping clean performance nearly unchanged.

\begin{table}[t]
\centering
\caption{Comparison between HSeCo, original DVCL, and tuned DVCL on ACM. Values are attack-average Micro-F1/Macro-F1 over PRBCD and HetePRBCD.}
\label{tab:lambda_compare}
\begin{tabular}{c|c|c}
\hline
Model & Attack Micro-F1 Avg & Attack Macro-F1 Avg \\
\hline
HSeCo & 0.8525 & 0.8488 \\
DVCL, $\lambda_{\mathrm{DVCL}}=0.1$ & 0.8665 & 0.8632 \\
DVCL, $\lambda_{\mathrm{DVCL}}=0.5$ & 0.8707 & 0.8680 \\
\hline
\end{tabular}
\end{table}

We then perform component ablation under the raw directed KNN setting with $k=20$. Table~\ref{tab:component_ablation} shows that the full DVCL model outperforms the variant without contrastive loss and both single-view variants. This confirms that the main gain comes from the combination of the topology view, the feature view, and cross-view contrastive alignment.

\begin{table}[t]
\centering
\caption{ACM component ablation under raw KNN feature view.}
\label{tab:component_ablation}
\begin{tabular}{c|c|c|c}
\hline
Variant & Clean Micro-F1 & Attack Micro-F1 Avg & Attack Macro-F1 Avg \\
\hline
Full DVCL & $0.8770\pm0.0071$ & 0.8710 & 0.8680 \\
DVCL w/o CL & $0.8743\pm0.0074$ & 0.8656 & 0.8626 \\
Topology only & $0.8523\pm0.0177$ & 0.8471 & 0.8415 \\
Feature only & $0.8539\pm0.0086$ & 0.8539 & 0.8512 \\
\hline
\end{tabular}
\end{table}

\subsection{Parameter Sensitivity}

Table~\ref{tab:param} reports sensitivity to $\lambda_{\mathrm{DVCL}}$ under the raw directed KNN setting. Increasing the contrastive weight improves the attack-average score, but larger weights also bring higher variance on some HetePRBCD settings. We therefore use $\lambda_{\mathrm{DVCL}}=0.5$ as the current main configuration and treat larger values as candidates for further tuning.

\begin{table}[t]
\centering
\caption{Sensitivity to $\lambda_{\mathrm{DVCL}}$ on ACM.}
\label{tab:param}
\begin{tabular}{c|c|c|c}
\hline
$\lambda_{\mathrm{DVCL}}$ & Clean Micro-F1 & Attack Micro-F1 Avg & Attack Macro-F1 Avg \\
\hline
0.5 & $0.8770\pm0.0071$ & 0.8710 & 0.8680 \\
1.0 & $0.8810\pm0.0049$ & 0.8723 & 0.8690 \\
2.0 & $0.8790\pm0.0085$ & 0.8734 & 0.8709 \\
\hline
\end{tabular}
\end{table}

\subsection{Supplementary Feature-View Upgrade}

The main DVCL configuration uses raw directed KNN with $k=20$. We also evaluate whether the feature view can be further improved by changing the KNN construction. These experiments are supplementary upgrades, not the main baseline comparison.

Table~\ref{tab:knn_supp} reports a lightweight KNN study. The raw $k=20$ feature view is the default setting. Heterogeneous-neighbor-enhanced KNN improves attack-average Micro-F1, suggesting that the feature view can be upgraded by incorporating direct heterogeneous neighborhood information.

\begin{table}[t]
\centering
\caption{Supplementary KNN feature-view study on ACM. Values are three-seed Micro-F1 averages.}
\label{tab:knn_supp}
\begin{tabular}{c|c|c|c}
\hline
Feature view & Clean & Attack Micro-F1 Avg & Attack Macro-F1 Avg \\
\hline
Raw KNN, $k=10$ & 0.8750 & 0.8719 & 0.8693 \\
Raw KNN, $k=20$ & 0.8752 & 0.8706 & 0.8675 \\
Raw KNN, $k=40$ & 0.8665 & 0.8677 & 0.8650 \\
Mutual raw KNN, $k=20$ & 0.8804 & 0.8739 & 0.8715 \\
Hetero-profile KNN, $k=20$ & 0.8807 & 0.8784 & 0.8766 \\
\hline
\end{tabular}
\end{table}

Later experiments further explore normalized heterogeneous profiles, mutual KNN, and relation-wise neighbor feature means. These variants can improve the feature view, but they should be treated as follow-up extensions of DVCL rather than the core method used in the main HSeCo comparison.

% ========== 5 Related Work ==========

